% PLP: original CCM periodic points and the complete common-label lattice carrier.
\section{Periodic adelic points with the full lattice coefficient}
\label{sec:cc-periodic-full-lattice}

{\raggedright
\paragraph{Human source and exact scope.}
Alain Connes, Caterina Consani, and Matilde Marcolli,
\href{https://arxiv.org/abs/math/0703392v1}{\emph{The Weil proof and the
geometry of the adeles class space}}, original author file
\nolinkurl{Yuri70v2.tex}, Section \emph{Thermodynamics and geometry of the
primes}, supply the periodic locus and the original one-zero-place
adeles. The author file was read at lines 2425--2944, including
\texttt{globalper}, \texttt{preorbitckak}, \texttt{classXiadeles},
\texttt{avwadele}, \texttt{Xiorbits}, \texttt{muscalehaar},
\texttt{classpointsperorb}, \texttt{kmsbeware}, and \texttt{XiKalgpts}.
The retained archive is version 1; the author's filename does not change
that version. Its SHA-256 is
\texttt{b1f694cb934dd4fe\allowbreak{}f829287f1a55dd8d\allowbreak{}e21c67a481e55245\allowbreak{}3c2f7e8fe0ae6fb5}.
The actual programme coefficient and sheaf maps are those proved in
\nolinkurl{ARBITRARY_SUPPORT_CC_COEFFICIENTS.tex}, ULC1--10 and ULC17,
and \nolinkurl{FULL_SUPPORT_SHEAF_COMPARISON.tex}, FSC1--9 and FSC33--34.
Those complete local derivations accompany this calculation. All new
carrier, orbit, and localization statements below are proved here.
No positive-energy, KMS, Hilbert-space positivity, or RH assertion is
deduced from a support label.
\par}
\subsection{The original carrier and the different local-product object}

Let $L$ be a nontrivial bounded distributive lattice, with bottom $0_L$
and top $1_L$. It need not be finite. For a commutative ring $R$ retain
the constrained carrier and its full operations
\begin{equation}\tag{PLP1}
\begin{gathered}
G_L(R)=\{(0,\lambda):\lambda\in L\}\cup
                  \{(a,1_L):a\in R\},\\
(a,\lambda)+(b,\mu)=(a+b,\lambda\vee\mu),\qquad
(a,\lambda)(b,\mu)=(ab,\lambda\wedge\mu),\\
\tau=(0,0_L),\qquad e=(0,1_L),\qquad {\bf1}=(1,1_L).
\end{gathered}
\end{equation}
The two descriptions of $(0,1_L)$ denote the same element.
The constraint is $a\ne0\Rightarrow\lambda=1_L$.
It is preserved by both displayed operations: if a product has
nonzero amplitude both factors have top label, and if a sum has
nonzero amplitude at least one summand has top label. Distributivity
follows from ring distributivity and lattice distributivity. In
particular $e^2=e$, $e+\tau=e$, and $e\ne\tau$.

Write $\Sigma=\{\infty\}\cup\{p:p\text{ rational prime}\}$ and
$\mathbb A=\mathbb A_{\mathbb Q}$. Define separately
\begin{equation}\tag{PLP2}
S_{\mathrm c}=G_L(\mathbb A),\qquad
S_{\mathrm{loc}}=
G_L(\mathbb R)\times\prod_p' G_L(\mathbb Q_p),
\quad\text{restricted relative to }G_L(\mathbb Z_p).
\end{equation}
The second expression is a newly declared semiring, not another name
for the first. More explicitly, its elements are tuples
$((a_v,\lambda_v))_{v\in\Sigma}$ for which $a=(a_v)\in\mathbb A$
and $a_v\ne0\Rightarrow\lambda_v=1_L$. There is no additional
restriction on the zero labels, since every $(0,\lambda_v)$ lies in
$G_L(\mathbb Z_p)$. The usual finite-exception condition for the
amplitudes is preserved under coordinate addition and multiplication;
PLP1 therefore proves that this restricted product is a semiring.

There is an exact comparison
\begin{equation}\tag{PLP3}
\Delta:S_{\mathrm c}\longrightarrow S_{\mathrm{loc}},
\qquad (a,\lambda)\longmapsto((a_v,\lambda))_v.
\end{equation}
It is an injective unital semiring homomorphism. Indeed each coordinate
respects PLP1; the amplitudes and the common label can be recovered from
the tuple. Its image is exactly the tuples with one constant label
$\lambda_v=\lambda$ for all $v$. If that common label is not top, the
constraint at each coordinate forces all amplitudes to vanish, which
proves the reverse inclusion without enlarging $G_L(\mathbb A)$.
The same map is linear for the diagonal $G_L(\mathbb Z)$ action.

Let $\rho:S_{\mathrm{loc}}\to\mathbb A$ and
$q:S_{\mathrm c}\to\mathbb A$ be amplitude projection. Both are
surjective semiring homomorphisms and $\rho\Delta=q$. For
$Z(a)=\{v:a_v=0\}$ the exact fibres are
\begin{equation}\tag{PLP4}
\rho^{-1}(a)=L^{Z(a)},\qquad
q^{-1}(a)=
\begin{cases}
\{(a,1_L)\},&a\ne0,\\
\{(0,\lambda):\lambda\in L\},&a=0.
\end{cases}
\end{equation}
The first identification inserts top labels at the coordinates outside
$Z(a)$. At $a=0$ the map induced by $\Delta$ is the diagonal
$L\hookrightarrow L^\Sigma$. At nonzero $a$ its image is the single
point with top label at every coordinate, including every zero
coordinate. These are the full fibres, not only selected examples.

There is no amplitude-preserving zero-preserving semiring
homomorphism $T:S_{\mathrm{loc}}\to S_{\mathrm c}$. To prove this,
fix a place $v$ and define $x$ to have entry $(1,1_L)$ at $v$ and
$\tau$ elsewhere; define $y$ to have entry $\tau$ at $v$ and
$(1,1_L)$ elsewhere. They belong to the restricted product, and
$xy=\tau_{\mathrm{loc}}$. Their amplitude adeles $a,b$ are nonzero,
$a+b=1$, and $ab=0$. Amplitude preservation forces
\begin{equation}\tag{PLP5}
T(x)=(a,1_L),\quad T(y)=(b,1_L),\quad
T(x)T(y)=e,\quad T(xy)=\tau.
\end{equation}
This contradicts $e\ne\tau$. In particular there is no
amplitude-preserving semiring retraction of $\Delta$. The injection
PLP3 and the fibre comparison PLP4 remain valid despite this obstruction.

\subsection{Actual support localization and its infinite-product defect}

Both carriers are algebras over $S=G_L(\mathbb Z)$ via diagonal
integers and common labels. Put $E=\Delta(e)$. Their first support
localizations are
\begin{equation}\tag{PLP6}
S_{\mathrm c}[e^{-1}]\simeq L,\qquad
S_{\mathrm{loc}}[E^{-1}]\simeq L^\Sigma.
\end{equation}
Multiplication by $e$ or $E$ erases the amplitudes and preserves
every existing label. An invertible idempotent is the unit; hence in
either localization each element equals that product. Conversely
equality between zero-label tuples is unchanged by multiplying by
$E$, so no two distinct tuples become equal at this stage.
Every label tuple is present already with all amplitudes zero. This
proves both isomorphisms, including injectivity. The map induced by
$\Delta$ is the diagonal $L\to L^\Sigma$.

Let $J$ be a proper prime ideal of the lattice $L$ and
$P_J=\{(0,\lambda):\lambda\in J\}\subset S$. Then
\begin{equation}\tag{PLP7}
\begin{gathered}
(S_{\mathrm c})_{P_J}=L_J,\qquad
(S_{\mathrm{loc}})_{P_J}=
\mathscr L_J:=
(L^\Sigma)[\operatorname{diag}(L\setminus J)^{-1}],\\
[\boldsymbol\lambda]=[\boldsymbol\mu]\text{ in }\mathscr L_J
\ \Longleftrightarrow\
\exists\nu\notin J\ \forall v\in\Sigma:
\lambda_v\wedge\nu=\mu_v\wedge\nu .
\end{gathered}
\end{equation}
First invert $e$, which is outside $P_J$. Every nonzero top-supported
integer then maps to the unit. The remaining denominators are exactly
the diagonal zero labels outside $J$. Products of finitely many such
denominators are again one diagonal $\nu\notin J$, because the
complement of a prime lattice ideal is closed under meets. Inverting
an idempotent identifies it with the unit, so every fraction has a
label-tuple numerator. The usual localization equality is consequently
precisely the displayed uniform-denominator relation. This also proves
the first formula with a single label, as in ULC6.

Coordinate localization gives a surjection
\begin{equation}\tag{PLP8}
\pi_J:\mathscr L_J\longrightarrow\prod_{v\in\Sigma}L_J,\qquad
[\boldsymbol\lambda]\longmapsto([\lambda_v]_J)_v.
\end{equation}
Surjectivity uses coordinate representatives, available by countable
choice. Equality after this map means
$\forall v\ \exists\nu_v\notin J:
\lambda_v\wedge\nu_v=\mu_v\wedge\nu_v$; equality before it requires
one common $\nu$ in PLP7. For a finite set of coordinates the meet
of the finitely many $\nu_v$ is outside $J$ and is a common
denominator. Infinite coordinates do not permit that argument.
There is also injectivity for every place set when $L\setminus J$
has a least element $\nu_0$: every coordinate equality then holds
after meeting with $\nu_0$. In particular this applies to finite
$L$, since the meet of all elements outside $J$ is still outside
$J$. This finite-lattice fact does not extend to arbitrary $L$.

Here is an explicit failure for the actual countable place set.
Take the bounded totally ordered lattice
\begin{equation}\tag{PLP9}
L=\{0_L,1_L,a_1,a_2,\ldots\},
\qquad 0_L<\cdots<a_3<a_2<a_1<1_L,\qquad J=\{0_L\}.
\end{equation}
Joins and meets are maximum and minimum; $J$ is prime since a meet
of two positive elements is positive. Enumerate the places $v_n$ and
set $\lambda_{v_n}=a_n$, $\mu_{v_n}=1_L$. At coordinate $n$, meeting
with $\nu_n=a_n$ proves equality in $L_J$. A uniform denominator
would have to be a positive element below every $a_n$, and there is
no such element in this displayed lattice. Thus $\pi_J$ is not
injective. The exact defect is the difference between the two
quantifier orders in PLP7--PLP8; replacing the left-hand side by
the product of stalks loses this defect.

The diagonal $L_J\to\mathscr L_J$ is still injective: equality of
two constant tuples by a common denominator is exactly equality
in $L_J$. For tuples which are top except at one fixed place $v$,
PLP7 reduces to
\begin{equation}\tag{PLP10}
[(1,\ldots,\lambda\text{ at }v,\ldots)]
=[(1,\ldots,\mu\text{ at }v,\ldots)]
\ \Longleftrightarrow\ [\lambda]_J=[\mu]_J.
\end{equation}
Indeed every other coordinate gives the tautology $\nu=\nu$.
The infinite-product defect therefore does not change this exact
one-place equality criterion.

\subsection{The sheaf comparison over the original split base}

Use the notation of ULC1--3:
$X=\operatorname{Spec}\mathbb Z$, $S_L=\operatorname{Spec}_{\rm lat}L$,
$Y=\operatorname{Spec}G_L(\mathbb Z)=S_L\sqcup X$,
and $W_U=S_L\cup U$ for nonempty arithmetic opens $U\subset X$.
Use the actual structure sheaf $\mathcal O_Y$, whose restriction
to $S_L$ is $\mathcal O_{S_L}$, with
$\mathcal O_{S_L}(D_L(\lambda))=\downarrow\lambda$.
The adelic point amplitude $\mathbb A$ is a rational vector space;
no complex scalar multiplication on every $\mathbb Q_p$ is asserted.
Its constant sheaf on $X$ is an $\mathcal O_X$-module using the
inclusion $\mathcal O_X(U)\subset\mathbb Q$.

The ULC construction on this module gives $\mathcal M_{\mathbb A}$.
Define the separate locally labelled sheaf $\mathcal N$ by
\begin{equation}\tag{PLP11}
\begin{array}{lll}
\mathcal M_{\mathbb A}(W_U)=S_{\mathrm c},&
\mathcal M_{\mathbb A}(V)=\mathcal O_{S_L}(V),& V\subset S_L,\\
\mathcal N(W_U)=S_{\mathrm{loc}},&
\mathcal N(V)=\displaystyle\prod_{v\in\Sigma}\mathcal O_{S_L}(V),
& V\subset S_L.
\end{array}
\end{equation}
Restrictions between arithmetic opens are identities.
Restriction to support retains the common label for the first sheaf,
and retains every coordinate label for the second, followed by the
actual support restrictions. The scalar action on arithmetic opens
is diagonal multiplication by $(r,\mu)\in
G_L(\mathcal O_X(U))$; on support opens it is coordinatewise meet
with the one scalar section of $\mathcal O_{S_L}$.

These formulas define sheaves of $\mathcal O_Y$-semimodules.
For a cover of $W_U$, its arithmetic members have nonempty pairwise
arithmetic intersections because $\operatorname{Spec}\mathbb Z$
is irreducible. Compatible sections on them are the same carrier
element. Each arithmetic member contains all of $S_L$, so that
compatibility forces every support-only member to be the specified
restriction. The sections consequently glue uniquely. Covers within
$S_L$ glue by the sheaf axiom for $\mathcal O_{S_L}$, coordinate
by coordinate for the product. The formulas for the scalar operations
commute with all these restrictions, which proves the assertion.

The map $\Delta$ extends to a sheaf monomorphism
\begin{equation}\tag{PLP12}
\mathcal M_{\mathbb A}\hookrightarrow\mathcal N;
\quad\text{on }W_U\text{ it is PLP3, on }V\subset S_L
\text{ it is the diagonal support map}.
\end{equation}
Both its arithmetic and support section maps are injective and
compatible with the restrictions just proved. At a support prime
the stalks and comparison are exactly
\begin{equation}\tag{PLP13}
(\mathcal M_{\mathbb A})_J=L_J,\qquad
\mathcal N_J=\mathscr L_J,\qquad
L_J\hookrightarrow\mathscr L_J
\xrightarrow{\pi_J}\prod_{v\in\Sigma}L_J .
\end{equation}
For the second equality, basic neighbourhoods $D_L(\nu)$,
$\nu\notin J$, are cofinal. Their sections are
$(\downarrow\nu)^\Sigma$ and the restrictions meet every coordinate
with the same smaller basic-open label. Every such section is the
restriction of its label tuple in $L^\Sigma$. Equality of two germs
therefore has one basic neighbourhood witnessing equality of all
coordinates at once, which is precisely PLP7. This directly proves
the stalk statement and the failure of commuting this infinite
product with the stalk functor in PLP9.

For each $v$ the same construction with constant rational module
$\mathbb Q_v$ gives $\mathcal M_{\mathbb Q_v}$. Coordinate
projection yields a commuting diagram of actual sheaf maps:
\begin{equation}\tag{PLP14}
\mathcal M_{\mathbb A}\longrightarrow\mathcal N
\longrightarrow\mathcal M_{\mathbb Q_v},\qquad
(a,\lambda)\longmapsto(a_v,\lambda).
\end{equation}
At arithmetic points the second arrow projects the $v$ entry. At
support points it is $\mathscr L_J\to L_J$, with composite
$L_J\to L_J$ the identity. Thus the original support map
$(a,\lambda)\mapsto[\lambda]_J$ of ULC8 and FSC34 is retained
exactly. Its target is not the unchanged adele amplitude.

\subsection{All labels above the classical periodic points}

Write $I=\mathbb A^\times$, $C=I/\mathbb Q^\times$. The units of
both carriers are exactly the top-supported ideles. For a unit,
the label of it and of its inverse must have meet $1_L$, so both
labels at every coordinate are top. The amplitudes must be mutual
inverses in the adele ring. Conversely an idele and its inverse
with top labels multiply to $\mathbf{1}$. Requiring every local
amplitude to be nonzero alone would be insufficient: its inverse
must also satisfy the restricted-product condition.

Multiplication by $j\in I$ acts on either carrier by
$(a,\lambda)\mapsto(ja,\lambda)$, coordinatewise in the local
carrier. These are additive semimodule automorphisms and commute
with $\Delta$ and amplitude projection. They are not asserted to
be ring automorphisms for pointwise multiplication; for example
$j(ab)$ need not equal $(ja)(jb)$. The quotient objects below
are orbit sets, not semiring quotients.

The maps induced on rational orbit sets are
\begin{equation}\tag{PLP15}
S_{\mathrm c}/\mathbb Q^\times
\ \lhook\joinrel\longrightarrow\
S_{\mathrm{loc}}/\mathbb Q^\times
\longrightarrow
\mathbb A/\mathbb Q^\times .
\end{equation}
The first remains injective: an equality between two images by one
rational multiplier is, by equivariance and injectivity of PLP3,
the same equality in the original carrier. The image consists of
constant-label orbits. The fibre over a nonzero amplitude orbit
$[a]$ is a singleton for the first carrier and $L^{Z(a)}$ for
the local carrier. A nonzero adele has trivial rational stabilizer,
since any one of its nonzero coordinates forces $q=1$; labels
are unchanged by multiplication. Hence the labels in this fibre
are independent of the representative. The fibres over $[0]$
are respectively $L$ and $L^\Sigma$.

For a place $v$, the source's exact adele is
\begin{equation}\tag{PLP16}
a^{(v)}_w=
\begin{cases}0,&w=v,\\1,&w\ne v.\end{cases}
\end{equation}
It is nonzero as an adele. Its only lift in $S_{\mathrm c}$ is
$(a^{(v)},1_L)$, and $\Delta$ assigns the supported zero
$e=(0,1_L)$ to its $v$ coordinate. In $S_{\mathrm{loc}}$ the
whole fibre is the explicitly different family
\begin{equation}\tag{PLP17}
a^{(v)}_\lambda:
\quad (0,\lambda)\text{ at }v,\quad (1,1_L)\text{ at }w\ne v,
\qquad\lambda\in L .
\end{equation}
Only $\lambda=1_L$ is in the image of the original carrier.
The member with $\lambda=0_L$ exists in the enlarged local carrier;
it is not a second lift inside $G_L(\mathbb A)$.
All original global zero elements $(0,\lambda)$ and all local
zero-amplitude tuples $((0,\lambda_v))_v$ are also present.
Every idele fixes them, including $e$ and $\tau$ as distinct points.
At a support prime the original lift of $a^{(v)}$ maps to $1$ in
$L_J$; the local lifts map to their one-place label tuples in
$\mathscr L_J$, with the exact equality relation PLP10.

For completeness, every nonzero amplitude $a$ has rational
stabilizer $\{1\}$ and raw idele stabilizer
\begin{equation}\tag{PLP18}
I_{Z(a)}=\{j\in I:j_w=1\text{ for every }w\notin Z(a)\}.
\end{equation}
At $[a]$, its $I$ stabilizer is $\mathbb Q^\times I_{Z(a)}$
and its $C$ stabilizer is the injective image of $I_{Z(a)}$.
Indeed $ja=qa$ forces $j_w=q$ at every nonzero coordinate,
and imposes no equation at zero coordinates. Moreover a rational
idele supported on $Z(a)$ equals $1$ at a nonzero coordinate,
so its rational value is $1$. These arguments do not depend on
the labels. Thus a nonzero orbit has a nonidentity $C$ stabilizer
exactly when at least one coordinate vanishes. This proves the
support-labelled version of the source's periodic-locus assertion
without selecting one preferred label. The all-zero cases have
the whole group as stabilizer.

In particular, if $j_v(h)$ denotes the idele equal to $h$ at $v$
and $1$ elsewhere, then for every lift in PLP16--PLP17,
\begin{equation}\tag{PLP19}
\begin{gathered}
\operatorname{Stab}_{\mathbb Q^\times}(a^{(v)}_\lambda)=\{1\},
\qquad
\operatorname{Stab}_{I}(a^{(v)}_\lambda)=j_v(\mathbb Q_v^\times),\\
\operatorname{Stab}_{I}([a^{(v)}_\lambda])
=\mathbb Q^\times j_v(\mathbb Q_v^\times),\qquad
\operatorname{Stab}_{C}([a^{(v)}_\lambda])
=\iota_v(\mathbb Q_v^\times),\\
O_v=C/\iota_v(\mathbb Q_v^\times).
\end{gathered}
\end{equation}
Here $\iota_v$ is injective by the rational-intersection argument.
For the original carrier the subscript $\lambda$ means its sole
value $1_L$. At the rational quotient every rational already acts
trivially; the first formula is a statement about the raw point.
The local enlarged carrier has disjoint label-indexed copies
$\{O_{v,\lambda}:\lambda\in L\}$, all with these stabilizers.
The original carrier has precisely $O_{v,1_L}$.

\subsection{The real flow, compact factors, and every orientation}

The canonical section of the idele class group gives
$C=\mathbb R_{>0}\times K$ with
$K=\prod_\ell\mathbb Z_\ell^\times$. To prove this directly, for
an idele $j$ multiply by the rational
$\operatorname{sgn}(j_\infty)\prod_p p^{-v_p(j_p)}$.
Only finitely many exponents are nonzero. Its real coordinate is
positive and all its finite coordinates are units. A rational
number which preserves these conditions has every finite valuation
zero and positive sign, hence is $1$. The representative is unique,
is compatible with multiplication, and its real coordinate is the
full idele norm $|j_\infty|\prod_p|j_p|_p$.
The section and its inverse are continuous: in the idele topology
the finite valuation tuple and real sign are locally constant, and
on such a neighbourhood the section is multiplication by one
fixed rational number followed by continuous coordinate maps.

For $h=p^n\epsilon\in\mathbb Q_p^\times$,
$\epsilon\in\mathbb Z_p^\times$, the exact image is
\begin{equation}\tag{PLP20}
\iota_p(p^n\epsilon)=
\left(p^{-n},\,u_p=\epsilon,\,
u_\ell=p^{-n}\quad(\ell\ne p)\right).
\end{equation}
This follows by multiplying $j_p(h)$ by the rational $p^{-n}$,
not by changing the norm convention. Put
$K^{(p)}=\prod_{\ell\ne p}\mathbb Z_\ell^\times$ and
$\ell_p=\log p$. Quotienting by the $p$-unit part of the stabilizer
and using the generator in PLP20 gives the homogeneous-space
presentation
\begin{equation}\tag{PLP21}
O_p=(\mathbb R\times K^{(p)})/
\bigl((x,u)\sim(x+\ell_p,pu)\bigr).
\end{equation}
Indeed the stabilizer generator in these coordinates is
$(-\ell_p,p^{-1})$, whose inverse gives the displayed relation.
Specify the forward real flow by $g_t=(e^t,1)\in C$.
It acts by $[x,u]\mapsto[x+t,u]$, so the return to the
$x=0$ transversal after positive time $\ell_p$ is exactly
\begin{equation}\tag{PLP22}
[\,\ell_p,u\,]=[\,0,p^{-1}u\,],\qquad
u\longmapsto p^{-1}u .
\end{equation}
Multiplication by $p$ would describe the opposite time or the
opposite return convention.

The full $O_p$ has no nonzero pure-real period. If
$(e^t,1)=\iota_p(p^n\epsilon)$, the $p$ coordinate forces
$\epsilon=1$, and any finite $\ell\ne p$ forces $p^{-n}=1$
as an element of $\mathbb Q_\ell$, whence $n=0$ and $t=0$.
Since $C$ is abelian this applies at every point. Its compact
quotient instead has
\begin{equation}\tag{PLP23}
O_p/K=\mathbb R_{>0}/p^{\mathbb Z}
=\mathbb R/(\log p)\mathbb Z ,
\end{equation}
with least positive forward period $\log p$. It is this compact
quotient which the original CCM discussion after \texttt{XiKalgpts}
identifies as the periodic real-flow orbit. This distinction holds
for every label copy, without changing any norm.

At the infinite place the exact formula is
\begin{equation}\tag{PLP24}
\iota_\infty(h)=
\left(|h|,(\operatorname{sgn}h)_p\right),\qquad
O_\infty=K/\{1,-1\},\qquad O_\infty/K=\{\mathrm{point}\}.
\end{equation}
Multiply the original idele by the rational sign of $h$ to
obtain this representative. The whole positive-real subgroup
lies in its stabilizer and fixes $O_\infty$ pointwise.
Equations PLP21--PLP24 give the quotient-group topology of each
homogeneous orbit. Its orbit map to the original
$\mathbb A/\mathbb Q^\times$ is continuous. We retain that
ambient quotient topology, as the source requires; the equations
do not replace the embedded union by a disjoint-union topology.
For arbitrary $L$, no additional topology on its element set is
postulated. Each label-preserving copy has the specified homogeneous
orbit parametrization; the base's support topology remains the
actual lattice spectrum from PLP11--PLP14.

\subsection{Exact normal coordinates, their Haar factors, and their labels}

Put $\mathbb A_v=\{a\in\mathbb A:a_v=0\}$ and let
$b_v:\mathbb Q_v\to\mathbb A$ insert a coordinate at $v$ and
zero elsewhere. With the original additive topology and measures,
\begin{equation}\tag{PLP25}
\mathbb A=\mathbb A_v\oplus b_v(\mathbb Q_v),\qquad
\mathbb A/\mathbb A_v=\mathbb Q_v,\qquad
a^{(v)}(z)=a^{(v)}+b_v(z).
\end{equation}
Projection and insertion are continuous and inverse to the displayed
product decomposition: all finite restricted-product conditions
are unchanged by altering one coordinate. The map
$z\mapsto[a^{(v)}(z)]$ into the rational orbit space is injective,
because a rational identification forces that rational to equal
$1$ at every coordinate $w\ne v$.
For $h\in\mathbb Q_v^\times$,
\begin{equation}\tag{PLP26}
j_v(h)a^{(v)}(z)=a^{(v)}(hz),\qquad
(1-h):z\longmapsto(1-h)z,\qquad
d_v((1-h)z)=|1-h|_v\,d_vz\quad(h\ne1).
\end{equation}
For $h\ne1$ the difference map has inverse multiplication by
$(1-h)^{-1}$. Its exact additive-Haar determinant is
$|1-h|_v$, using real Lebesgue measure or
$\operatorname{vol}(\mathbb Z_p)=1$, respectively.

One can retain the whole transverse family as an associated
action space
\begin{equation}\tag{PLP27}
\mathcal T_v=
C\times_{\iota_v(\mathbb Q_v^\times)}\mathbb Q_v,\qquad
(g\iota_v(h),z)\sim(g,hz),\qquad
[g,z]\longmapsto[g\,a^{(v)}(z)] .
\end{equation}
The last expression uses an idele representative of $g$ and is
independent of that representative. The relation is respected by
PLP26. At $z=0$ the map is the bijection with $O_v$ of PLP19.
At $z\ne0$, $a^{(v)}(z)=j_v(z)$ is an idele, and the map is
$[g,z]\mapsto g\iota_v(z)\in C$. This is a bijection: equality
of the products implies exactly the displayed associated-space
relation. No point of this part can equal a zero-coordinate
point. Continuity follows from the continuous map
$I\times\mathbb Q_v\to\mathbb A/\mathbb Q^\times$,
$(j,z)\mapsto[j a^{(v)}(z)]$. It is unchanged by $j\mapsto qj$.
The quotient $I\to C$ is open, since the inverse image of the
image of an open subset is its union of rational translates.
Its product with the identity of $\mathbb Q_v$ is open and is
a quotient map. Hence the displayed map descends continuously
to $C\times\mathbb Q_v$, then through the associated relation.
Corestriction to its image with the ambient subspace topology
is continuous. Thus PLP27 is a continuous bijection onto
$C\cup O_v$ inside the adelic orbit space. No assertion that
its inverse is continuous in the ambient orbit topology is
required here.

At a finite prime write the full representative as
$g=(e^x,\epsilon,u)$ with $\epsilon\in\mathbb Z_p^\times$ and
$u\in K^{(p)}$. Since $\iota_p(\epsilon)=(1,\epsilon,1)$,
the associated relation gives exactly
\[
 [(e^x,\epsilon,u),z_{\rm old}]
 =[(e^x,1,u),z_{\rm new}],\qquad
 z_{\rm new}=\epsilon z_{\rm old},\qquad
 z_{\rm old}=\epsilon^{-1}z_{\rm new},\qquad
 |\epsilon|_p=1 .
\]
Thus the original $p$-unit is recorded in the normal coordinate,
with its exact Haar factor retained. Writing $z=z_{\rm new}$,
the remaining representative coordinates $(x,u,z)$ obey
\begin{equation}\tag{PLP28}
(x+\ell_p,pu,z)\sim(x,u,p^{-1}z),\qquad
(u,z)\longmapsto(p^{-1}u,p^{-1}z)
\quad\text{after time }+\ell_p .
\end{equation}
To check the sign, changing $(x,u)$ by
$(+\ell_p,p)$ multiplies its class by $\iota_p(p^{-1})$.
The associated relation then multiplies $z$ by $p^{-1}$.
For every positive integer $n$ the two opposite-direction
normal determinants are
\begin{equation}\tag{PLP29}
|1-p^n|_p=1,\qquad |1-p^{-n}|_p=p^n,\qquad
|1-p^n|_p^{-1}=1,\qquad |1-p^{-n}|_p^{-1}=p^{-n}.
\end{equation}
The first follows because $1-p^n$ is a $p$-adic unit; the
second follows by writing $1-p^{-n}=(p^n-1)/p^n$.
The forward return is the second case. Negative iterates exchange
these cases; the zeroth iterate has zero difference map and is
excluded from the displayed inverse-determinant formula.

For comparison with a trace on test functions retain the inverse
in the pullback action:
$\vartheta(h)f(z)=f(h^{-1}z)$. Its additive distribution kernel is
$\delta(y-h^{-1}z)$. At $h\ne1$ the transverse diagonal is
\begin{equation}\tag{PLP30}
\int_{\mathbb Q_v}\delta((1-h^{-1})z)\,d_vz
=|1-h^{-1}|_v^{-1}
=\frac{|h|_v}{|1-h|_v}.
\end{equation}
The first equality is the defining change of variables for
the Dirac distribution under an invertible scalar; it can
equally be checked against a compactly supported test that is
one near zero. For a compactly supported coefficient $b$ on
$\mathbb Q_v^\times$ whose support avoids $1$, multiplicative
Haar measure is invariant under inversion, so
\begin{equation}\tag{PLP31}
\int_{\mathbb Q_v^\times}
\frac{b(h)}{|1-h^{-1}|_v}\,d^\times h
=
\int_{\mathbb Q_v^\times}
\frac{b(h^{-1})}{|1-h|_v}\,d^\times h .
\end{equation}
This proves the relation between the inverse test argument and the
normal denominator away from the singular identity. It is not a
replacement for the source's regularization there or for its
global trace theorem.

The original hyperplane Haar factor is also retained in full:
for $r\in\mathbb Q^\times$,
\begin{equation}\tag{PLP32}
d_{\mathbb A_v}(ra)
=\left(\prod_{w\ne v}|r|_w\right)d_{\mathbb A_v}a
=|r|_v^{-1}d_{\mathbb A_v}a,\qquad
|r|_\infty\prod_p|r|_p=1 .
\end{equation}
Only finitely many finite-place factors differ from $1$.
Writing $r=\pm\prod_p p^{n_p}$ gives the last product exactly
as $(\prod_p p^{n_p})(\prod_p p^{-n_p})=1$.
The product decomposition PLP25 then proves the first formula.
This is the factor in the human source's \texttt{muscalehaar};
no measure on support labels has been inserted.

The supported normal-coordinate map is the surjective unital
semiring map
\begin{equation}\tag{PLP33}
\nu_v^\uparrow:G_L(\mathbb A)\longrightarrow G_L(\mathbb Q_v),
\qquad (a,\lambda)\longmapsto(a_v,\lambda).
\end{equation}
A lift is $(b_v(z),\lambda)$; when $z\ne0$ the original constraint
already forces $\lambda=1_L$. The insertion is additive and
multiplicative but is not unital: its unit is the adele idempotent
supported at $v$. The exact equality relation for PLP33 is
$(a,\lambda)\sim(c,\mu)$ if and only if
$\lambda=\mu$ and $a-c\in\mathbb A_v$.
This is a semiring congruence, since coordinate projection
respects addition and multiplication. Thus its same-label
quotient is exactly $G_L(\mathbb Q_v)$; a quotient which also
identifies different zero labels is a different construction.
The map is PLP14 on arithmetic stalks. On support stalks it is
the identity of $L_J$. It is rationally equivariant, not
rationally invariant, and therefore is not asserted to give
an unquotiented local-field coordinate on all of
$\mathbb A/\mathbb Q^\times$. The specific transversal PLP25
is what supplies that actual coordinate.

In the original common-label carrier the affine map
$(z,\mu)\mapsto(a^{(v)},1_L)+(b_v(z),\mu)$ is
$(a^{(v)}(z),1_L)$ for every admissible $(z,\mu)$.
Thus it forgets the direction label when $z=0$; it cannot produce
the missing local points $a^{(v)}_\lambda$ for $\lambda<1_L$.
In $S_{\mathrm{loc}}$, insert $(z,\mu)$ at $v$ and $\tau$
elsewhere and add it to the actual point $a^{(v)}_\lambda$.
The resulting map is
\begin{equation}\tag{PLP34}
\gamma_{v,\lambda}(z,\mu)_w=
\begin{cases}
(z,\lambda\vee\mu),&w=v,\\
(1,1_L),&w\ne v.
\end{cases}
\end{equation}
Its precise fibre relation is equality of $z$ and of
$\lambda\vee\mu$. The $\tau$ direction leaves the base point
fixed; the $e$ direction moves it to $a^{(v)}_{1_L}$.
Every nonzero $z$ has $\mu=1_L$ and therefore top support.
The local stabilizer acts by $z\mapsto hz$ and leaves these
labels unchanged, so this affine map is equivariant.

At $\lambda=0_L$, the map PLP34 is injective on the entire
normal carrier. It yields the exact enlarged associated family
\begin{equation}\tag{PLP35}
\widetilde{\mathcal T}_v=
C\times_{\iota_v(\mathbb Q_v^\times)}G_L(\mathbb Q_v)
\longrightarrow S_{\mathrm{loc}}/\mathbb Q^\times,
\quad
[g,(z,\mu)]\longmapsto[g\,\gamma_{v,0_L}(z,\mu)].
\end{equation}
It is a bijection onto $C\cup\bigcup_{\mu\in L}O_{v,\mu}$.
For $z\ne0$ this is the nonzero part proved in PLP27, with its
forced top label. For $z=0$ the label $\mu$ is fixed by every
group action, and the quotient by the local group is precisely
$O_{v,\mu}$; different labels cannot be identified. The map
which retains amplitudes gives a surjection
$\widetilde{\mathcal T}_v\to\mathcal T_v$, is bijective off
the zero section, and has fibre $L$ over each point of the zero
section. Its injection back as action spaces selects
$(z,1_L)$ for every $z$, and identifies the original family
PLP27 with $C\cup O_{v,1_L}$. The section $z\mapsto(z,1_L)$
sends $0$ to $e$, not $\tau$, so it is not falsely claimed to
be a zero-preserving semiring map. Equations PLP34--PLP35
compute the exact new family defined by the failed retraction
in PLP5.

Finally the supported version of the transverse difference
retains its value even at the identity:
\begin{equation}\tag{PLP36}
(z,\mu)+(-h,1_L)(z,\mu)=((1-h)z,\mu).
\end{equation}
For $h\ne1$ this is an invertible top-supported scalar map of
$G_L(\mathbb Q_v)$. For $h=1$ it is multiplication by $e$,
namely $(z,\mu)\mapsto(0,\mu)$, and not the map constantly
equal to $\tau$. On the support stalk $L_J$, both nonzero
top-supported scalars and $e$ act as the identity. This is
compatible with the exact localization PLP6--PLP14; ordinary
subtraction on an idempotent lattice has not been assumed.
The additive-Haar determinants PLP26--PLP32 concern the
original local-field amplitudes and do not assign an
unproved determinant or multiplicity to the label set.

\subsection{The receiving maps in the original CC coefficient sheaves}

The adelic points and the CC test-function coefficients are related
by evaluation, not by identifying their amplitude spaces. Write
$V=\mathcal S(\mathbb A)$ and let $\mathcal H_k=\mathcal O_k$
be the divisor-indexed CC function sheaf in CDV and ULC.
Its restrictions are inclusions of subspaces of $V$.
For each $a\in\mathbb A$, evaluation $\xi\mapsto\xi(a)$
therefore defines a morphism
$\mathcal H_k\to\underline{\mathbb C}_X$ of
$\mathcal O_X$-modules. The constant target is a sheaf because
every nonempty arithmetic open is irreducible.
The ULC construction supplies the exact lift
\begin{equation}\tag{PLP37}
\operatorname{ev}_a^\uparrow:
\mathcal M_k\longrightarrow\mathcal M_{\mathbb C},
\qquad
(\xi,\lambda)\longmapsto(\xi(a),\lambda)
\text{ on }W_U,\qquad
L_J\longrightarrow L_J\text{ the identity at support}.
\end{equation}
The target lies in the original carrier: a nonzero evaluated
amplitude implies a nonzero original function and hence top
label. Evaluation is additive and rational-linear; the label
join and meet formulas are unchanged. Restriction to support
on both sides is $(\xi,\lambda)\mapsto[\lambda]_J$, proving
compatibility with every ULC8/FSC34 map.

Retain the original rational dilation convention
$\alpha_r\xi(x)=\xi(r^{-1}x)$, $d(r)_p=v_p(r)$.
ULC17 supplies the exact map
$\alpha_r:\mathcal M_k\to\mathcal M_{k+d(r)}$.
Point multiplication and function dilation then satisfy
\begin{equation}\tag{PLP38}
\operatorname{ev}_{ra}^\uparrow\circ\alpha_r
=\operatorname{ev}_a^\uparrow,\qquad
(\alpha_r\xi)(ra)=\xi(r^{-1}ra)=\xi(a).
\end{equation}
The equation holds on each arithmetic open and on each support
stalk, where all maps preserve the original label.
It connects the rational orbit computation to the actual
divisor-indexed sheaves with no change of amplitude or support.
Evaluation through $\rho$ on the enlarged local carrier has
the same value on all $a^{(v)}_\lambda$, because their original
adele is identical. The carrier $\mathcal N$ records the extra
labels; ordinary evaluation through $\mathbb A$ does not.
No evaluation homomorphism of a crossed-product algebra is
asserted by these linear sheaf maps.

\paragraph{Precisely retained conclusion.}
The original arbitrary-$L$ carrier contains the CCM classical
orbit union with the unique common top support at each nonzero
adele. Its exact injection into the independently labelled
restricted product has fibres, support-stalk maps, and an explicit
non-retraction defect given in PLP3--PLP14. The latter construction
has the whole $L$ family over each one-zero-place orbit, with its
proved associated normal family PLP35. The original prime orbit,
the $\log p$ period after the compact quotient, both return
orientations, the local normal factor, the original Haar scaling,
and the receiving coefficient-sheaf maps have all been retained.
These constructions do not prove any new positivity assertion.
